\chapter{Results}
\label{ch:results}

This chapter reports what was observed. It is organised by research question: the demand and
workload characteristics that condition everything else (\Cref{sec:res-demand}), the capacity
results answering RQ1 (\Cref{sec:res-capacity}), the policy comparison answering RQ2
(\Cref{sec:res-policy}), the bottleneck and adaptive-capacity results answering RQ3
(\Cref{sec:res-bottleneck}), and the forecast-based planning results answering RQ4
(\Cref{sec:res-future}). Interpretation is deferred to \Cref{ch:discussion}; the aim here is to
state the evidence and the patterns it contains.

All retrospective results come from a single run covering twelve months, sixteen workforce
candidates per month and three policies --- 576 simulation runs in total --- under the economic
assumptions recorded in \Cref{sec:economics}. All forecast results come from a separate
December run under \emph{different} economic assumptions. As established in
\Cref{sec:economics}, monetary quantities are never compared between the two.

\section{Demand and workload characteristics}
\label{sec:res-demand}

\Cref{tab:demand-profile} summarises the year. Volume varies by a factor of nearly five between
the summer trough (\num{8400} orders in July and August) and the December peak
(\num{40800}). The urgent share moves with volume rather than against it, from
\SI{7.9}{\percent} in July to \SI{25.4}{\percent} in December, so the peak months combine the
greatest volume with the highest concentration of the tightly constrained class.

\begin{table}[htbp]
\centering
\caption[Monthly demand profile]{Monthly demand profile of the dataset used for the
retrospective experiments. Workload units are means per order.}
\label{tab:demand-profile}
\small
\begin{tabular}{lrrrrrr}
\toprule
\textbf{Month} & \textbf{Orders} & \textbf{Urgent} & \textbf{Mean} &
\multicolumn{3}{c}{\textbf{Mean workload units}} \\
\cmidrule(lr){5-7}
 & & \textbf{(\%)} & \textbf{items} & \textbf{Picking} & \textbf{Packing} & \textbf{Dispatch} \\
\midrule
January   & \num{31200} & 20.0 & 3.59 & 4.09 & 2.70 & 1.20 \\
February  & \num{26400} & 18.0 & 3.59 & 4.07 & 2.68 & 1.19 \\
March     & \num{16800} & 12.1 & 2.80 & 3.16 & 2.36 & 1.16 \\
April     & \num{16800} & 12.4 & 2.80 & 3.17 & 2.38 & 1.16 \\
May       & \num{10800} &  9.0 & 2.17 & 2.45 & 2.14 & 1.14 \\
June      & \num{9600}  &  8.7 & 2.20 & 2.47 & 2.13 & 1.14 \\
July      & \num{8400}  &  7.9 & 2.21 & 2.48 & 2.13 & 1.14 \\
August    & \num{8400}  &  8.1 & 2.22 & 2.49 & 2.15 & 1.14 \\
September & \num{16800} & 12.2 & 2.99 & 3.38 & 2.44 & 1.16 \\
October   & \num{20400} & 15.2 & 3.02 & 3.43 & 2.48 & 1.18 \\
November  & \num{33600} & 22.3 & 3.99 & 4.55 & 2.86 & 1.21 \\
December  & \num{40800} & 25.4 & 4.37 & 5.01 & 3.03 & 1.22 \\
\midrule
\textbf{Year} & \num{240000} & 17.2 & 3.20 & 3.39 & 2.46 & 1.17 \\
\bottomrule
\end{tabular}
\end{table}

The three stages are loaded very differently by the same order stream. Averaged over the year,
one order carries \num{3.39} workload units at picking, \num{2.46} at packing and \num{1.17} at
dispatch, a ratio of \(1 : 0.72 : 0.34\). The ratio is not constant: between May and December
the picking workload per order rises by \SI{104}{\percent} while dispatch workload rises by
only \SI{7}{\percent}. Seasonal demand growth therefore falls very unevenly across the three
stages, concentrating on picking.

Two mechanisms produce this concentration, and they compound. The first is direct: picking
workload scales most strongly with item count, and the mean order grows from \num{2.17} items in
May to \num{4.37} in December. The second is compositional: the urgent share doubles over the
same period, and urgent orders are drawn from a more demanding mix --- \SI{35}{\percent} fragile
against \SI{20}{\percent} for normal orders, and \SI{30}{\percent} high-complexity against
\SI{15}{\percent}. The peak months therefore do not simply contain more orders; they contain
larger, more complex, more urgent orders, each of which loads picking disproportionately.

The consequence for capacity planning is that the three stages cannot be scaled together. Moving
from the July trough to the December peak requires the total order volume to be multiplied by
almost five, but the picking workload per order also doubles, so picking capacity must grow far
faster than dispatch capacity. This is visible in the recommendations of
\Cref{tab:recommendations}: the picking-to-dispatch worker ratio in the recommended
configurations widens from \(3:2\) in July to \(22:5\) in December.

\section{Capacity planning: analytical estimate against simulation (RQ1)}
\label{sec:res-capacity}

\subsection{The recommendations}

\Cref{tab:recommendations} reports, for each month, the analytical estimate and the
configuration that the framework selected after simulating all sixteen candidates under all
three policies. The framework's recommendation rule is the lexicographic one of
\Cref{sec:selection-rule}: cheapest among the configurations meeting both service floors, with a
least-violation fallback only if none does. The fallback was not reached in any month --- a
feasible configuration existed in all twelve. The table itself is tabulated from the export's
lowest-total-cost column, which is computed without a feasibility filter; in this run it
coincides with the lexicographic recommendation in every month, configuration for configuration
(\Cref{sec:app-rules}).

\begin{table}[htbp]
\centering
\caption[Monthly workforce recommendations]{Monthly recommendations from the retrospective
workflow. ``Analytical'' is the total FTE from \Cref{eq:capacity-estimate}; the remaining
columns describe the configuration selected after simulation.}
\label{tab:recommendations}
\small
\begin{tabular}{lrrlrrrr}
\toprule
\textbf{Month} & \textbf{Orders} & \textbf{Analytical} & \textbf{Configuration} &
\textbf{Policy} & \textbf{FTE} & \textbf{Urgent} & \textbf{Cost} \\
 & & \textbf{FTE} & & & & \textbf{SLA (\%)} & \textbf{(EUR)} \\
\midrule
January   & \num{31200} & 31 & \rgm{s13\_8\_4} & RL-3 & 25 & 99.9 & \num{74085} \\
February  & \num{26400} & 27 & \rgm{s13\_8\_4} & RL-3 & 25 & 99.9 & \num{67065} \\
March     & \num{16800} & 15 & \rgm{s643}      & UF   & 13 & 100.0 & \num{31955} \\
April     & \num{16800} & 15 & \rgm{s643}      & UF   & 13 & 100.0 & \num{32440} \\
May       & \num{10800} &  9 & \rgm{s332}      & RL-3 &  8 & 99.7 & \num{20695} \\
June      & \num{9600}  &  9 & \rgm{s322}      & RL-3 &  7 & 99.9 & \num{17475} \\
July      & \num{8400}  &  8 & \rgm{s322}      & RL-3 &  7 & 100.0 & \num{16955} \\
August    & \num{8400}  &  8 & \rgm{s322}      & UF   &  7 & 100.0 & \num{16965} \\
September & \num{16800} & 16 & \rgm{s743}      & UF   & 14 & 100.0 & \num{34390} \\
October   & \num{20400} & 19 & \rgm{s863}      & UF   & 17 & 100.0 & \num{41655} \\
November  & \num{33600} & 36 & \rgm{s16\_9\_5} & RL-3 & 30 & 99.9 & \num{77905} \\
December  & \num{40800} & 47 & \rgm{s22\_11\_5}& RL-3 & 38 & 99.2 & \num{115040} \\
\bottomrule
\end{tabular}
\end{table}

\subsection{The analytical estimate systematically over-sizes the workforce}

The clearest result of the capacity analysis is visible in \Cref{fig:analytical}: in
\emph{every one} of the twelve months, the configuration selected after simulation is smaller
than the analytical estimate. The gap is never negative and never zero. It is modest in the low
season --- one FTE in May, July and August --- and substantial in the peak: six FTE in January,
six in November and nine in December.

\begin{figure}[htbp]
\centering
\includegraphics[width=\textwidth]{f3_analytical_vs_recommended.pdf}
\caption[Analytical estimate against simulated recommendation]{Analytical capacity anchor
against the workforce recommended after simulation. The annotation on each pair gives the
difference. The estimate exceeds the simulated recommendation in all twelve months.}
\label{fig:analytical}
\end{figure}

The magnitude is material. December's nine-FTE difference corresponds to
\(9 \times \num{2400} = \eur{21600}\) of monthly labour cost under that run's economic
assumptions. Against an analytical estimate of 47 FTE, a simulated recommendation of 38 is
\(9/47 \approx \SI{19}{\percent}\) below the estimate; expressed instead as a share of the
recommended workforce the same nine FTE are \(9/38 \approx \SI{24}{\percent}\). The first of
these --- the difference as a proportion of the analytical estimate --- is the form used
consistently throughout this thesis. Aggregated across the
peak months alone, the difference between sizing by formula and sizing by simulation is
measured in tens of thousands of euros per year.

Two mechanisms account for the gap, and they work in the same direction. First, the analytical
estimate deliberately targets \SI{85}{\percent} utilisation, which builds in a safety margin.
Second, and more interestingly, the estimate has no way of knowing that a sequencing policy can
protect the urgent class. It sizes capacity as though all orders were interchangeable, whereas
the simulated evaluation discovers that a smaller workforce combined with an urgency-aware
sequencing rule satisfies both service floors. The estimate cannot represent that substitution
between capacity and operating discipline, and consequently over-provides capacity.

\subsection{Where the over-estimation falls}
\label{sec:res-stage-gap}

Decomposing the gap by stage shows that it is not distributed evenly.
\Cref{tab:stage-gap} gives the per-stage difference between the analytical estimate and the
recommended configuration.

\begin{table}[htbp]
\centering
\caption[Per-stage difference between estimate and recommendation]{Analytical estimate against
recommended configuration, by stage. Values are picking / packing / dispatch FTE.}
\label{tab:stage-gap}
\small
\begin{tabular}{lccc}
\toprule
\textbf{Month} & \textbf{Analytical} & \textbf{Recommended} & \textbf{Difference} \\
\midrule
January   & 16 / 10 / 5 & 13 / 8 / 4 & \(-3\) / \(-2\) / \(-1\) \\
February  & 14 / 8 / 5  & 13 / 8 / 4 & \(-1\) / \(0\) / \(-1\) \\
March     & 7 / 5 / 3   & 6 / 4 / 3  & \(-1\) / \(-1\) / \(0\) \\
April     & 7 / 5 / 3   & 6 / 4 / 3  & \(-1\) / \(-1\) / \(0\) \\
May       & 4 / 3 / 2   & 3 / 3 / 2  & \(-1\) / \(0\) / \(0\) \\
June      & 4 / 3 / 2   & 3 / 2 / 2  & \(-1\) / \(-1\) / \(0\) \\
July      & 3 / 3 / 2   & 3 / 2 / 2  & \(0\) / \(-1\) / \(0\) \\
August    & 3 / 3 / 2   & 3 / 2 / 2  & \(0\) / \(-1\) / \(0\) \\
September & 8 / 5 / 3   & 7 / 4 / 3  & \(-1\) / \(-1\) / \(0\) \\
October   & 9 / 6 / 4   & 8 / 6 / 3  & \(-1\) / \(0\) / \(-1\) \\
November  & 19 / 11 / 6 & 16 / 9 / 5 & \(-3\) / \(-2\) / \(-1\) \\
December  & 26 / 14 / 7 & 22 / 11 / 5 & \(-4\) / \(-3\) / \(-2\) \\
\midrule
\textbf{Total} & & & \(\mathbf{-17}\) / \(\mathbf{-13}\) / \(\mathbf{-6}\) \\
\bottomrule
\end{tabular}
\end{table}

Summed across the year, the estimate over-provides by 17 picking FTE, 13 packing FTE and 6
dispatch FTE. In absolute terms picking absorbs the largest correction, which is unsurprising
since it is the largest stage. In \emph{relative} terms, however, packing is over-provided most
consistently: it is reduced in eight of the twelve months, including the summer months where
picking is not reduced at all.

The reason is that the analytical estimate rounds each stage up independently. A stage whose
computed requirement is only slightly above an integer receives a full additional worker, and
the smaller the stage, the larger that rounding is as a proportion of its capacity. Packing,
sitting between a large picking stage and a small dispatch stage, is most exposed to this
effect. In July and August the estimate calls for three packing FTE where two suffice --- a
third of the stage's capacity, obtained purely from rounding.

This is worth noting because it identifies a second, quite separate source of the gap. The
substitution between capacity and sequencing discipline explains the large corrections in the
peak months; independent per-stage rounding explains the persistent one-worker corrections in
the quiet ones. Both push in the same direction, and neither is visible without simulating.

\section{Sequencing policies (RQ2)}
\label{sec:res-policy}

Every comparison in this section is made under common random numbers, so the three policies face
identical order arrivals and identical service times for every order at every stage. Where two
policies are compared at the same workforce configuration, the only difference between the runs
is the order in which waiting work was selected. Differences reported below therefore reflect
the sequencing decision itself and not variation in the demand or service realisation. The
evidence that this synchronisation holds is visible in the results: at a given configuration the
three policies produce stage utilisations that agree to within a tenth of a percentage point,
which would not occur if they were drawing different service times.

Two complementary views are reported. The equal-workforce view isolates the effect of sequencing
by holding capacity constant. The minimum-feasible-workforce view
(\Cref{sec:res-min-feasible}) asks what each policy costs in capacity terms, which is the form
in which the comparison enters a budget discussion. The two answer different questions and are
kept separate throughout.

\subsection{Feasibility}

The single most decisive result concerns how often each policy can satisfy the service floors
at all. \Cref{tab:feasibility} and \Cref{fig:feasibility} report, for each month, the number of
the sixteen candidate configurations that met both floors under each policy.

\begin{table}[htbp]
\centering
\caption[Feasible configurations by month and policy]{Number of the sixteen candidate
configurations meeting both service floors, by month and policy.}
\label{tab:feasibility}
\small
\begin{tabular}{lrrr@{\hskip 2em}lrrr}
\toprule
\textbf{Month} & \textbf{FIFO} & \textbf{UF} & \textbf{RL-3} &
\textbf{Month} & \textbf{FIFO} & \textbf{UF} & \textbf{RL-3} \\
\midrule
January   & \textbf{0} & 16 & 16 & July      &  9 & 11 & 11 \\
February  & \textbf{0} & 15 & 15 & August    &  9 & 11 & 11 \\
March     & 11 & 14 & 14 & September & 11 & 16 & 16 \\
April     & 11 & 14 & 14 & October   & 14 & 16 & 16 \\
May       &  8 & 11 & 11 & November  & \textbf{0} & 16 & 16 \\
June      & 11 & 11 & 11 & December  & \textbf{0} & 15 & 16 \\
\midrule
\multicolumn{8}{l}{\textbf{Total across all 192 runs per policy:}
FIFO \textbf{84}, Urgent-First \textbf{166}, RL-3 \textbf{167}} \\
\bottomrule
\end{tabular}
\end{table}

\begin{figure}[htbp]
\centering
\includegraphics[width=\textwidth]{f4_feasibility_by_month.pdf}
\caption[Feasible configurations by month and policy]{Configurations meeting both service
floors, by month and policy. In the four campaign months FIFO satisfies the floors at none of
the sixteen candidates.}
\label{fig:feasibility}
\end{figure}

Two patterns stand out. First, aggregated over all 192 runs per policy, FIFO is feasible in
\num{84} while Urgent-First and RL-3 are feasible in \num{166} and \num{167} respectively ---
roughly double.

Second, and more striking, FIFO's failures are not spread evenly. Its \num{108} infeasible
configurations occur in every month of the year --- \num{44} of them across the eight
non-campaign months, so it would be wrong to say the failures fall entirely in peak --- but what
happens
\emph{only} in peak is total infeasibility. In January, February, November and December, the
four campaign months and precisely the months with the highest volume and the highest urgent
share, FIFO satisfies the service floors at \emph{none} of the sixteen candidate workforces,
whereas in every other month at least eight of the sixteen are feasible. Within the evaluated
range, adding capacity does not rescue it in those four months. The reason is one of discipline
rather than of capacity: a first-in-first-out queue gives an urgent order no precedence over the
large volume of normal orders ahead of it, and with a \SI{240}{\minute} threshold against a
\SI{1440}{\minute} one, the urgent class fails first and fails hardest. A workforce large enough
to keep every queue short would eventually satisfy both floors under FIFO too; none of the
candidates a planner would consider for these months is that large.

\subsection{The class-level trade-off}

\Cref{fig:tradeoff} plots urgent against normal service attainment for all 576 runs. The
structure of the plot is the finding. FIFO observations are dispersed down the vertical axis:
normal attainment is frequently high --- often above \SI{90}{\percent} --- while urgent
attainment falls as low as \SI{10}{\percent}. Urgent-First and RL-3 observations are compressed
into a band along the top, with urgent attainment near \SI{100}{\percent} almost everywhere,
and vary instead along the horizontal axis according to how much normal service they retain.

\begin{figure}[htbp]
\centering
\includegraphics[width=0.82\textwidth]{f5_class_sla_tradeoff.pdf}
\caption[Class-level service trade-off]{Urgent against normal service attainment for all 576
runs. The shaded region satisfies both floors. FIFO sacrifices the urgent class; the two
urgency-aware policies differ mainly in how much normal service they retain.}
\label{fig:tradeoff}
\end{figure}

This shows that the operative question is not whether to protect the urgent class --- both
urgency-aware policies do that almost completely --- but what protecting it costs the normal
class. That is where Urgent-First and RL-3 separate, and it is the mechanism behind the
December result below.

\subsection{Three representative capacity regimes}

Aggregate counts conceal how strongly the answer depends on the capacity regime.
\Cref{tab:cases} reports three configurations chosen to span the range, and
\Cref{fig:cases} shows two of them.

\begin{table}[htbp]
\centering
\caption[Representative capacity regimes]{Three configurations spanning the capacity range.
All three policies face identical stochastic realisations under common random numbers.}
\label{tab:cases}
\small
\begin{tabular}{llrlrrrlr}
\toprule
\textbf{Month} & \textbf{Config.} & \textbf{FTE} & \textbf{Policy} &
\textbf{Urgent} & \textbf{Normal} & \textbf{Total} & \textbf{Feasible} & \textbf{Cost} \\
 & & & & \textbf{(\%)} & \textbf{(\%)} & \textbf{(\%)} & & \textbf{(EUR)} \\
\midrule
\multirow{3}{*}{October} & \multirow{3}{*}{\rgm{s753}} & \multirow{3}{*}{15}
  & FIFO & 24.5 & 91.1 & 81.0 & no  & \num{90460} \\
 & & & UF   & 100.0 & 84.2 & 86.6 & yes & \num{49635} \\
 & & & RL-3 & 99.7 & 89.4 & 90.9 & yes & \textbf{\num{45400}} \\
\midrule
\multirow{3}{*}{October} & \multirow{3}{*}{\rgm{s11\_7\_5}} & \multirow{3}{*}{23}
  & FIFO & 99.9 & 100.0 & 100.0 & yes & \num{55275} \\
 & & & UF   & 100.0 & 99.9 & 100.0 & yes & \textbf{\num{55245}} \\
 & & & RL-3 & 99.9 & 100.0 & 100.0 & yes & \num{55260} \\
\midrule
\multirow{3}{*}{December} & \multirow{3}{*}{\rgm{s22\_11\_5}} & \multirow{3}{*}{38}
  & FIFO & 19.7 & 89.7 & 71.9 & no  & \num{273140} \\
 & & & UF   & 99.3 & 73.8 & 80.2 & no  & \num{132570} \\
 & & & RL-3 & 99.2 & 85.4 & 88.9 & yes & \textbf{\num{115040}} \\
\bottomrule
\end{tabular}
\end{table}

\begin{figure}[htbp]
\centering
\includegraphics[width=\textwidth]{f7_representative_cases.pdf}
\caption[Policy behaviour under two capacity regimes]{Class-level attainment under two October
configurations. The dashed line marks the urgent floor (\SI{95}{\percent}) and the dotted line
the normal floor (\SI{80}{\percent}). At \rgm{s753} the policies differ dramatically; at
\rgm{s11\_7\_5} they are indistinguishable.}
\label{fig:cases}
\end{figure}

\paragraph{Under-capacity (October, \rgm{s753}, 15 FTE).} The policies diverge sharply. FIFO
delivers \SI{24.5}{\percent} urgent attainment and is infeasible. Urgent-First achieves perfect
urgent attainment but retains only \SI{84.2}{\percent} normal service. RL-3 gives up a fraction
of a percentage point of urgent attainment (\SI{99.7}{\percent}) and recovers more than five
percentage points of normal service (\SI{89.4}{\percent}), and is the cheapest of the three at
\eur{45400} against Urgent-First's \eur{49635}.

\paragraph{Well-resourced (October, \rgm{s11\_7\_5}, 23 FTE).} The policies are practically
indistinguishable. All three exceed \SI{99.9}{\percent} on both classes, and their total costs
lie within \eur{30} of one another --- a spread of \SI{0.05}{\percent}. When capacity is
sufficient, queues rarely form, so there is seldom a choice to make and the sequencing rule has
almost nothing to act on.

\paragraph{The peak (December, \rgm{s22\_11\_5}, 38 FTE).} This configuration produces the
sharpest separation observed anywhere in the study. FIFO collapses to \SI{19.7}{\percent}
urgent attainment. Urgent-First protects the urgent class successfully
(\SI{99.3}{\percent}) but, in doing so, drives normal attainment down to
\SI{73.8}{\percent} --- \emph{below the \SI{80}{\percent} floor} --- and is therefore
infeasible. RL-3 achieves essentially the same urgent attainment (\SI{99.2}{\percent}) while
holding normal service at \SI{85.4}{\percent}, and is the only feasible policy at this
configuration. It is also the cheapest, at \eur{115040} against Urgent-First's \eur{132570}, a
difference of \SI{13.2}{\percent}.

This is the strongest single result for the learned policy: at the peak month's recommended
workforce, the implemented RL-3 policy reaches a feasible outcome that the implemented
Urgent-First baseline does not reach at all, and does so at lower cost.

How far that difference can be attributed to \emph{learning} requires care, and the answer is
that it cannot be attributed to learning alone. The two implementations differ on two dimensions
at once, as \Cref{sec:urgent-first,sec:rl3} set out: RL-3 chooses which class to serve as a
function of state, and it maintains first-in-first-out order within each class; Urgent-First
applies fixed class precedence, and its within-class order follows the internal structure of a
binary heap. Both dimensions plausibly contribute to the observed normal-class gap, and this
experiment cannot decompose the contribution of each. The result reported here is a difference
between the two policies \emph{as implemented}, not an isolated estimate of what dynamic class
selection is worth. \Cref{sec:disc-uf-ordering} examines the evidence on the second mechanism,
and \Cref{sec:future-work} identifies the variant that would separate them.

\Cref{fig:cost-sla} places this configuration among December's sixteen candidates.

\begin{figure}[htbp]
\centering
\includegraphics[width=0.85\textwidth]{f6_cost_vs_sla.pdf}
\caption[December cost against service]{Total cost against total service attainment for all
sixteen December candidates under each policy. FIFO occupies a separate, expensive and
low-service region.}
\label{fig:cost-sla}
\end{figure}

\subsection{The systematic comparison over all matched configurations}
\label{sec:res-matched}

The three configurations of \Cref{tab:cases} were chosen to illustrate. Because every (month,
workforce) pair was evaluated under all three policies on identical stochastic realisations,
the grid also supports a systematic comparison: it consists of \textbf{192 matched triples},
in each of which the only difference between the three runs is the sequencing decision.
\Cref{tab:matched} summarises the pairwise differences across all of them.

\begin{table}[htbp]
\centering
\caption[Pairwise policy differences over 192 matched configurations]{Pairwise policy
differences across all 192 matched (month, workforce) configurations. Service differences are
in percentage points; cost differences in euros under the retrospective run's assumptions.}
\label{tab:matched}
\small
\begin{tabular}{lrrrrr}
\toprule
\textbf{Comparison} & \textbf{Cheaper} & \textbf{Mean cost} & \textbf{Urgent} &
\textbf{Mean urgent} & \textbf{Mean normal} \\
 & \textbf{in} & \textbf{difference} & \textbf{better in} & \textbf{difference} &
 \textbf{difference} \\
\midrule
RL-3 vs Urgent-First & 94 / 192  & \(-575\)     & 0 / 192   & \(-0.62\) & \(-0.70\) \\
RL-3 vs FIFO         & 169 / 192 & \(-22{,}638\) & 169 / 192 & \(+32.59\) & \(-1.53\) \\
Urgent-First vs FIFO & 173 / 192 & \(-22{,}063\) & 169 / 192 & \(+33.21\) & \(-0.83\) \\
\bottomrule
\end{tabular}
\end{table}

Three observations follow, and the second and third qualify the headline results reported above.

\paragraph{Against FIFO the advantage is unambiguous.} Both urgency-aware policies improve
urgent-class attainment by more than 32 percentage points on average and reduce total cost by
roughly \eur{22000} per configuration. The cost of that improvement to the normal class is
small --- between \num{0.83} and \num{1.53} percentage points on average.

\paragraph{Against Urgent-First, RL-3 never wins on the urgent class.} In \emph{none} of the
192 configurations does RL-3 achieve higher urgent-class attainment than strict priority, and
on average it is \num{0.62} percentage points lower. This is exactly what the priority rule is
designed to guarantee, and it confirms that RL-3's contribution is not better protection of the
urgent class. Its advantage, where it exists, comes from what it gives back to the normal class.

\paragraph{The overall averages conceal a sharp regional split.} Taken over all 192
configurations, RL-3's mean normal-class difference against Urgent-First is
\emph{negative} (\num{-0.70} points) even though it is better in 125 of them --- meaning its
losses are far larger than its gains. Splitting by feasibility resolves this completely, as
\Cref{tab:matched-split} shows.

\begin{table}[htbp]
\centering
\caption[Matched comparison split by feasibility]{RL-3 against Urgent-First, split by whether
either policy produced a feasible plan at that configuration.}
\label{tab:matched-split}
\small
\begin{tabular}{lrrrr}
\toprule
\textbf{Region} & \textbf{Configs} & \textbf{Mean normal} & \textbf{Worst normal} &
\textbf{Mean cost} \\
 & & \textbf{difference (pts)} & \textbf{difference (pts)} & \textbf{difference} \\
\midrule
At least one policy feasible & 167 & \(+1.84\)  & \(-0.08\)  & \(-1{,}991\) \\
Neither policy feasible      & 25  & \(-17.69\) & \(-33.69\) & \(+8{,}880\) \\
\bottomrule
\end{tabular}
\end{table}

\textbf{Every one of RL-3's material normal-class losses occurs where neither policy produced a
feasible plan.} In the 167 configurations where a viable plan existed, RL-3's normal-class
attainment was better by \num{1.84} percentage points on average and its worst case was a loss
of \num{0.08} points --- effectively never worse. In the 25 configurations where both policies
failed, RL-3 was worse in all of them, by \num{17.69} points on average, and cost \eur{8880}
more.

Those 25 configurations are severely under-resourced: backlog shares of 27 to 35 per cent and
picking utilisation averaging \SI{93.8}{\percent}. They are found in the low-demand months at
workforces well below the recommended level --- May \rgm{s232}, for example, where FIFO attains
\SI{4.0}{\percent} urgent service, Urgent-First \SI{100}{\percent} urgent but only
\SI{57.4}{\percent} normal, and RL-3 \SI{94.1}{\percent} urgent and \SI{23.7}{\percent} normal.
None of the three is remotely acceptable. The honest reading is that RL-3 degrades less
gracefully than strict priority when capacity is grossly insufficient, but that this occurs only
in a region where no configuration would be adopted.

\subsection{Policy divergence increases with capacity pressure}
\label{sec:res-pressure-bands}

The matched design also allows the regime dependence claimed in \Cref{sec:res-policy} to be
measured rather than illustrated. Grouping the 192 configurations by picking utilisation gives
\Cref{tab:pressure-bands} and the left panel of \Cref{fig:equal-workforce}.

The statistic in the first column of that table is the \emph{absolute} cost difference
\(\lvert C_{\text{RL-3}} - C_{\text{UF}}\rvert\). It measures how far apart the two policies
land, not which of them lands lower --- a distinction that turns out to matter here, and which
the signed comparison below addresses separately.

\begin{table}[htbp]
\centering
\caption[Policy divergence by capacity-pressure band]{Divergence between RL-3 and Urgent-First
at equal workforce, grouped by picking utilisation. ``Absolute gap'' is
\(\lvert C_{\text{RL-3}} - C_{\text{UF}}\rvert\), a measure of separation; ``signed'' is
\(C_{\text{UF}} - C_{\text{RL-3}}\), positive where RL-3 is cheaper, reported separately for
configurations at which at least one policy produced a feasible plan and for those at which
neither did.}
\label{tab:pressure-bands}
\small
\setlength{\tabcolsep}{4.5pt}
\begin{tabular}{lrrrrr}
\toprule
& & \multicolumn{2}{c}{\textbf{Absolute gap (EUR)}} &
\multicolumn{2}{c}{\textbf{Mean signed gap (EUR)}} \\
\cmidrule(lr){3-4}\cmidrule(lr){5-6}
\textbf{Picking utilisation} & \textbf{Configs} & \textbf{mean} & \textbf{per 1{,}000}
& \textbf{feasible} & \textbf{neither} \\
 & & & \textbf{orders} & \textbf{exists} & \textbf{feasible} \\
\midrule
below \SI{70}{\percent}      & 33 & 537        & 55  & $-0.5$ (31)   & $-\num{8608}$ (2) \\
\SIrange{70}{80}{\percent}   & 47 & \num{1484} & 96  & $+858$ (45)   & $-\num{14925}$ (2) \\
\SIrange{80}{90}{\percent}   & 69 & \num{3004} & 106 & $+\num{2988}$ (69) & --- \\
\SIrange{90}{95}{\percent}   & 11 & \num{5563} & 283 & $+\num{4611}$ (8)  & $-\num{7987}$ (3) \\
above \SI{95}{\percent}      & 32 & \num{6320} & 545 & $+\num{3630}$ (14) & $-\num{8388}$ (18) \\
\bottomrule
\end{tabular}
\end{table}

\begin{figure}[htbp]
\centering
\includegraphics[width=\textwidth]{f12_equal_workforce.pdf}
\caption[Equal-workforce comparison across all configurations]{Left: the absolute cost
difference between the two urgency-aware policies grows monotonically across the observed
capacity-pressure bands. Right: RL-3's normal-class advantage over Urgent-First holds wherever a
feasible plan exists and reverses where none does.}
\label{fig:equal-workforce}
\end{figure}

The mean absolute gap rises monotonically across the five observed bands, from \eur{537} below
\SI{70}{\percent} utilisation to \eur{6320} above \SI{95}{\percent} --- a factor of nearly
\textbf{twelve}. The relationship is not driven by a handful of extreme cases: it holds for the
band means, each computed over between 11 and 69 configurations.

Two qualifications keep this statistic honest. The first is that absolute euro gaps depend on
how much there is to be late about: order volume, urgency mix and total exposure to late
penalties all vary across months, and the high-utilisation bands are populated mainly by
campaign months. Normalising per \num{1000} orders addresses that directly, and the monotone
pattern survives it --- \eur{55} per \num{1000} orders below \SI{70}{\percent} utilisation
against \eur{545} above \SI{95}{\percent}. The separation is therefore not merely a volume
artefact.

The second is that separation is not the same as advantage. The signed columns of
\Cref{tab:pressure-bands} show that where at least one policy produced a feasible plan, the
difference does favour RL-3 and does grow with pressure, from indistinguishable below
\SI{70}{\percent} to roughly \eur{3000}--\eur{4600} in the top three bands. Where neither policy
reached feasibility, the sign reverses sharply and Urgent-First is the cheaper of the two by
between \eur{8000} and \eur{15000} on average. Since \num{18} of the \num{32} configurations
above \SI{95}{\percent} utilisation fall in that second group, the growth in the absolute gap in
the top band is driven as much by the reversal as by the advantage. The right panel of
\Cref{fig:equal-workforce} shows the same split directly.

The defensible reading is therefore that \emph{divergence} between the two urgency-aware
policies increases across the observed utilisation bands, and that within the feasible region
that divergence favours RL-3. Utilisation is used here as an observed grouping variable, not as
an established cause: the bands coincide with the campaign months, with high urgent shares and
with lean configurations, and this design cannot separate those. This is the quantitative form
of the claim made from the two October configurations in \Cref{tab:cases}. Sequencing policy is
close to a free choice when capacity is comfortable and a consequential one when it is not, and
the transition is gradual rather than abrupt.

\subsection{How much workforce each policy needs to become feasible}
\label{sec:res-min-feasible}

The comparisons above hold the workforce fixed, which isolates the sequencing effect. The
complementary question is what each policy costs in \emph{capacity}: how many full-time
equivalents it requires before it satisfies both service floors at all.
\Cref{tab:min-feasible} reports the smallest feasible configuration found for each policy.

\begin{table}[htbp]
\centering
\caption[Smallest feasible workforce by policy]{Smallest total workforce at which each policy
satisfied both service floors. A dash indicates that no candidate examined was feasible. The
final column gives the range of total FTE spanned by that month's sixteen candidates.}
\label{tab:min-feasible}
\small
\begin{tabular}{lrrrl}
\toprule
\textbf{Month} & \textbf{FIFO} & \textbf{Urgent-First} & \textbf{RL-3} &
\textbf{Candidate range} \\
\midrule
January   & --- & 25 & 25 & 25--37 \\
February  & --- & 25 & 25 & 21--33 \\
March     & 13 & 13 & 13 & 11--19 \\
April     & 13 & 13 & 13 & 11--19 \\
May       &  9 &  7 &  7 & 6--12 \\
June      &  7 &  7 &  7 & 6--12 \\
July      &  7 &  6 &  6 & 5--11 \\
August    &  7 &  6 &  6 & 5--11 \\
September & 14 & 12 & 12 & 12--20 \\
October   & 17 & 15 & 15 & 15--23 \\
November  & --- & 30 & 30 & 30--42 \\
December  & --- & \textbf{45} & \textbf{38} & 38--56 \\
\bottomrule
\end{tabular}
\end{table}

Where FIFO is feasible at all, it generally requires more capacity than the urgency-aware
policies to get there --- two additional FTE in May, September and October, one in July and
August. Over most of the year the two urgency-aware policies reach feasibility at exactly the
same workforce, which reinforces the earlier finding that they are close substitutes outside the
peak.

December is the exception, and it is a substantial one. Urgent-First does not reach feasibility
until 45 FTE, whereas RL-3 reaches it at 38 --- a difference of \textbf{seven full-time
equivalents}, or \eur{16800} per month in labour cost alone under that run's assumptions. The
mechanism is the one identified in \Cref{tab:cases}: at workforces between 38 and 44 FTE,
Urgent-First protects the urgent class successfully but pushes the normal class below its
\SI{80}{\percent} floor, and the only way it can lift the normal class back above the floor is
to buy enough additional capacity that both classes are served comfortably.

Two caveats bound this table. First, the search is over sixteen candidates per month, so a
figure equal to the lower end of the candidate range --- January, February, November and
December for the urgency-aware policies --- means only that the smallest configuration
\emph{examined} was already feasible, not that a smaller one would not have been. The December
comparison is unaffected by this, since both policies were evaluated over the same range and
Urgent-First failed across the lower part of it. Second, the entries are drawn from a single
demand realisation per month, so a marginal figure could shift under a different realisation.

\subsection{Where the cost difference actually comes from}
\label{sec:res-cost-decomposition}

Because the three policies are compared at identical workforce configurations, the labour
component of total cost is by construction the same for all three. Every euro of difference
between policies is therefore attributable to the late-order penalty term.
\Cref{tab:cost-split} makes this explicit at each month's recommended configuration.

\begin{table}[htbp]
\centering
\caption[Cost decomposition at the recommended configuration]{Cost decomposition at each
month's recommended configuration. Labour cost is identical across policies because the
workforce is the same; the entire difference lies in late-order penalties.}
\label{tab:cost-split}
\small
\begin{tabular}{llrrrr}
\toprule
\textbf{Month} & \textbf{Config.} & \textbf{Labour} & \multicolumn{3}{c}{\textbf{Late-order penalties (EUR)}} \\
\cmidrule(lr){4-6}
 & & \textbf{(EUR)} & \textbf{FIFO} & \textbf{Urgent-First} & \textbf{RL-3} \\
\midrule
January   & \rgm{s13\_8\_4}  & \num{60000} & \num{99340}  & \num{24955} & \textbf{\num{14085}} \\
February  & \rgm{s13\_8\_4}  & \num{60000} & \num{54620}  & \num{11385} & \textbf{\num{7065}} \\
March     & \rgm{s643}       & \num{31200} & \num{1065}   & \textbf{\num{755}} & \num{840} \\
April     & \rgm{s643}       & \num{31200} & \num{1500}   & \textbf{\num{1240}} & \num{1260} \\
May       & \rgm{s332}       & \num{19200} & \num{2425}   & \num{2145}  & \textbf{\num{1495}} \\
June      & \rgm{s322}       & \num{16800} & \num{825}    & \num{725}   & \textbf{\num{675}} \\
July      & \rgm{s322}       & \num{16800} & \num{175}    & \num{160}   & \textbf{\num{155}} \\
August    & \rgm{s322}       & \num{16800} & \num{270}    & \textbf{\num{165}} & \num{205} \\
September & \rgm{s743}       & \num{33600} & \num{1025}   & \textbf{\num{790}} & \num{820} \\
October   & \rgm{s863}       & \num{40800} & \num{1190}   & \textbf{\num{855}} & \num{920} \\
November  & \rgm{s16\_9\_5}  & \num{72000} & \num{117390} & \num{23890} & \textbf{\num{5905}} \\
December  & \rgm{s22\_11\_5} & \num{91200} & \num{181940} & \num{41370} & \textbf{\num{23840}} \\
\bottomrule
\end{tabular}
\end{table}

The table separates the year cleanly into two regimes. In the eight low- and shoulder-season
months, penalties are small in absolute terms --- under \eur{2500} against labour costs of
\eur{16800} to \eur{40800} --- and the three policies differ by amounts that are immaterial
next to the labour bill. In the four campaign months, penalties become the dominant cost
component and the spread between policies is enormous. In November, FIFO incurs \eur{117390} of
penalties against RL-3's \eur{5905}, a reduction of \SI{95}{\percent}; even against
Urgent-First's \eur{23890}, RL-3 removes a further \SI{75}{\percent}.

The practical reading is that the sequencing policy is an instrument for controlling the
\emph{penalty} component of cost, and that this component is negligible for most of the year and
dominant in the peak. A policy decision that appears to have almost no financial consequence in
July has a five-figure consequence in November, using exactly the same workforce.

\subsection{Backlog is nearly invariant to policy}
\label{sec:res-backlog}

A further observation clarifies what sequencing does and does not achieve. At each recommended
configuration, the share of orders left unfinished at the end of the operating horizon is almost
identical across the three policies: \SI{8.59}{\percent}, \SI{8.97}{\percent} and
\SI{8.95}{\percent} in January for FIFO, Urgent-First and RL-3 respectively, and
\SI{10.42}{\percent}, \SI{11.09}{\percent} and \SI{11.10}{\percent} in December. Stage
utilisation is likewise identical to within a tenth of a percentage point.

This is the expected consequence of a non-pre-emptive discipline. Reordering a queue does not
create worker-minutes; it changes chiefly \emph{which} work is completed and in what order. The
invariance observed is close but not exact --- the January figures differ across policies by up
to \num{0.38} percentage points, and the December figures by up to \num{0.68} --- which is what
one would expect with heterogeneous service times and a finite cut-off, where ordering can move
a small number of completions across the horizon boundary. The claim supported by these
experiments is therefore the empirical one: at the configurations tested, aggregate throughput
was very nearly a property of the workforce rather than of the policy.

The finding is worth stating because it delimits what sequencing can be expected to deliver. A
policy cannot create capacity, and no sequencing rule will clear a backlog that the workforce is
too small to process. What a policy determines is the allocation of the resulting lateness
across the two service classes --- which, when those classes carry different commitments and
different penalties, is precisely the decision that governs both feasibility and cost.

\subsection{Which policy is selected}

Across the twelve months, the recommended configuration used RL-3 in seven months (January,
February, May, June, July, November and December) and Urgent-First in five (March, April,
August, September and October). FIFO was selected in none. Because a feasible configuration
existed in every month, this split is the same under the lexicographic recommendation rule and
under the descriptive lowest-cost statistic (\Cref{sec:app-rules}).

The split is not arbitrary. RL-3 is selected in the two winter peaks, the pre-Christmas ramp
and the deep summer trough, while Urgent-First is selected in the moderate shoulder months. In
the shoulder months the margins are extremely small --- August differs by \eur{10}
(\eur{16955} against \eur{16965}) --- so the practical reading is that the two urgency-aware
policies are close substitutes over much of the year, and separate decisively only when the
normal-class floor comes under pressure.

\section{Bottleneck diagnosis and the marginal FTE (RQ3)}
\label{sec:res-bottleneck}

\subsection{Where pressure concentrates}

The diagnostic identified \textbf{picking as the primary pressure location in all twelve
months}, with pressure scores between \num{0.915} and \num{0.995}
(\Cref{tab:bottleneck}, \Cref{fig:pressure}). The margin over the second-ranked stage is large:
in December, picking scores \num{0.929} against packing's \num{0.307}.

\begin{table}[htbp]
\centering
\caption[Primary bottleneck by month]{Primary pressure location by month, with the components
of its score. Picking ranks first in every month.}
\label{tab:bottleneck}
\small
\begin{tabular}{lrrrr}
\toprule
\textbf{Month} & \textbf{Pressure} & \textbf{Utilisation} & \textbf{p95 wait} &
\textbf{Late-order wait} \\
 & \textbf{score} & \textbf{(\%)} & \textbf{(min)} & \textbf{share (\%)} \\
\midrule
January   & 0.983 & 95.8 & 993.8  & 99.8 \\
February  & 0.935 & 83.8 & 849.1  & 99.9 \\
March     & 0.986 & 96.7 & 236.0  & 99.5 \\
April     & 0.986 & 96.4 & 245.7  & 100.0 \\
May       & 0.995 & 98.8 & 235.3  & 100.0 \\
June      & 0.960 & 90.1 & 104.3  & 100.0 \\
July      & 0.915 & 79.5 & 71.9   & 98.4 \\
August    & 0.918 & 80.0 & 173.5  & 98.9 \\
September & 0.952 & 88.0 & 201.8  & 99.7 \\
October   & 0.978 & 94.5 & 216.8  & 100.0 \\
November  & 0.990 & 97.5 & 1083.2 & 100.0 \\
December  & 0.930 & 90.5 & 682.2  & 83.7 \\
\bottomrule
\end{tabular}
\end{table}

\begin{figure}[htbp]
\centering
\includegraphics[width=\textwidth]{f8_stage_pressure.pdf}
\caption[Stage Pressure by stage and month]{Upper panel: pressure score of each stage by month.
Lower panel: decomposition of the primary stage's score into its four weighted components.}
\label{fig:pressure}
\end{figure}

The late-order waiting share is the most striking column: between \SI{83.7}{\percent} and
\SI{100}{\percent} of all waiting time accumulated by late orders occurs at picking. Waiting in
this system is almost entirely a picking phenomenon; packing and dispatch contribute
negligibly, even though their utilisation is often substantial. The decomposition in the lower
panel of \Cref{fig:pressure} makes the reason clear: picking saturates the relative components
(p95 wait and queue length) in every month, so its score is driven by all four signals
simultaneously.

\subsection{The most-utilised stage is not always the identified bottleneck}
\label{sec:res-util-vs-pressure}

An instructive detail is that the diagnostic does not simply select the busiest stage. In three
of the twelve months the highest-utilised stage was \emph{not} the one ranked first, as
\Cref{tab:util-vs-pressure} shows.

\begin{table}[htbp]
\centering
\caption[Utilisation against pressure ranking]{Months in which the most-utilised stage was not
ranked as the primary pressure location.}
\label{tab:util-vs-pressure}
\small
\begin{tabular}{llrrrl}
\toprule
\textbf{Month} & \textbf{Stage} & \textbf{Utilisation} & \textbf{p95 wait} &
\textbf{Late-order} & \textbf{Pressure} \\
 & & \textbf{(\%)} & \textbf{(min)} & \textbf{wait share (\%)} & \textbf{rank} \\
\midrule
\multirow{2}{*}{June}      & Picking & 90.1 & 104.3 & 100.0 & \textbf{1} \\
                           & Packing & 98.1 & 84.4  & 0.0   & 2 \\
\midrule
\multirow{2}{*}{September} & Picking & 88.0 & 201.8 & 99.7  & \textbf{1} \\
                           & Packing & 96.0 & 117.8 & 0.3   & 2 \\
\midrule
\multirow{2}{*}{December}  & Picking  & 90.5 & 682.2 & 83.7 & \textbf{1} \\
                           & Dispatch & 98.6 & 426.4 & 0.0  & 2 \\
\bottomrule
\end{tabular}
\end{table}

In June, packing ran at \SI{98.1}{\percent} utilisation against picking's \SI{90.1}{\percent},
and yet picking was ranked first with a pressure score of \num{0.960} against packing's
\num{0.692}. The reason is visible in the components: picking accounted for
\SI{100}{\percent} of the waiting time accumulated by late orders and carried both the longest
p95 wait and the largest queue, so it saturated three of the four score components, while
packing's advantage on utilisation contributed only \(0.40 \times 0.981 = 0.393\).

This behaviour follows directly from how the score is weighted, and whether it is desirable
depends on what the diagnostic is for. A stage can be almost fully occupied without generating
waiting, if work arrives at it in a steady trickle; conversely a stage can be less occupied and
still be where orders pile up, if arrivals to it are bursty. Since service failures are caused
by waiting rather than by utilisation as such, ranking by accumulated waiting is defensible for
the purpose of explaining lateness. It does mean, however, that the diagnostic answers ``where
does waiting accumulate?'' rather than ``which stage is closest to its capacity limit?'', and
those are different questions that a planner might reasonably ask on different occasions.

\subsection{A bottleneck is not a hiring instruction}

The adaptive capacity search was triggered in nine of the twelve months. The three exceptions
--- February, July and August --- are explained exactly by the trigger condition of
\Cref{sec:adaptive}: their recommended configurations were feasible and their picking
utilisation fell below the \SI{85}{\percent} threshold, at \SI{83.8}{\percent},
\SI{79.5}{\percent} and \SI{80.0}{\percent} respectively (\Cref{tab:bottleneck}). Where a plan
is comfortably feasible at moderate utilisation, the framework does not spend simulation effort
asking whether to add capacity.

In each of the nine months where it did run, the search constructed a candidate with one
additional picking worker --- the identified bottleneck --- re-simulated it under all three
policies, and compared the outcome against the incumbent. Only one stage was ever tested,
because the margin between the first- and second-ranked stages never fell below the
\num{0.05} threshold that would have prompted a second candidate; the smallest observed margin
was \num{0.268}, in June. \textbf{In all nine months the additional worker was rejected.}
\Cref{tab:adaptive} and \Cref{fig:marginal} report the comparison.

\begin{table}[htbp]
\centering
\caption[Adaptive capacity search outcomes]{Every candidate tested by the adaptive search. Each
adds one picking FTE to that month's recommended configuration.}
\label{tab:adaptive}
\small
\begin{tabular}{llrrrl}
\toprule
\textbf{Month} & \textbf{Candidate} & \textbf{Added labour} & \textbf{Penalty} &
\textbf{Net} & \textbf{Outcome} \\
 & & \textbf{(EUR)} & \textbf{saved (EUR)} & \textbf{(EUR)} & \\
\midrule
January   & \rgm{s14\_8\_4}  & \num{2400} & \num{595}  & \num{1805} & rejected \\
March     & \rgm{s743}       & \num{2400} & \num{100}  & \num{2300} & rejected \\
April     & \rgm{s743}       & \num{2400} & \num{110}  & \num{2290} & rejected \\
May       & \rgm{s432}       & \num{2400} & \num{1230} & \num{1170} & rejected \\
June      & \rgm{s422}       & \num{2400} & \num{20}   & \num{2380} & rejected \\
September & \rgm{s843}       & \num{2400} & \num{25}   & \num{2375} & rejected \\
October   & \rgm{s963}       & \num{2400} & \num{95}   & \num{2305} & rejected \\
November  & \rgm{s17\_9\_5}  & \num{2400} & \num{2025} & \num{375}  & rejected \\
December  & \rgm{s23\_11\_5} & \num{2400} & \num{20}   & \num{2380} & rejected \\
\bottomrule
\end{tabular}
\end{table}

\begin{figure}[htbp]
\centering
\includegraphics[width=\textwidth]{f9_marginal_fte.pdf}
\caption[The marginal FTE decision]{Cost of adding one FTE at the identified bottleneck against
the late-order penalties it avoids. The additional worker was rejected in every month tested.}
\label{fig:marginal}
\end{figure}

The economics are mostly not close. An additional FTE costs \eur{2400} per month, and in six of
the nine months it avoided less than \eur{130} of penalties --- in June and December, only
\eur{20}. Two months are closer: May recovers \eur{1230} and November \eur{2025}, the latter
within \SI{16}{\percent} of breaking even.

The December case makes the underlying arithmetic explicit. At that month's recommended
configuration, one FTE costs \eur{2400}; avoiding that cost through reduced penalties would
require preventing \num{120} late urgent orders, or \num{480} late normal orders, or about
\num{456} orders in the observed mix. The observed average penalty per late order was
EUR~5.27. The additional capacity simply does not prevent enough failures to pay for
itself, notwithstanding that picking is unambiguously the most pressured stage.

\section{Forecast-based planning (RQ4)}
\label{sec:res-future}

The forecast workflow was exercised on December. Starting from an expected volume of
\num{40800} orders and the standard uncertainty setting, it generated three stochastic
replications and evaluated candidates in two stages: all sixteen candidates screened under all
three policies on the first replication, then the four strongest validated across all three.

The recommendation was \rgm{s25\_14\_7} --- 25 picking, 14 packing and 7 dispatch FTE ---
operated under RL-3, with expected total cost \eur{158103}, urgent attainment
\SI{99.9}{\percent} and normal attainment \SI{91.9}{\percent}. The configuration met both
service floors in all three replications.

\begin{figure}[htbp]
\centering
\includegraphics[width=0.85\textwidth]{f11_future_planning.pdf}
\caption[Forecast planning outcomes for December]{December candidates under forecast-based
planning: mean cost against the share of replications meeting both service floors. With three
replications that share takes only the values \(0\), \(\tfrac{1}{3}\), \(\tfrac{2}{3}\) and
\(1\). Circles were validated on three replications, crosses screened on one.}
\label{fig:future}
\end{figure}

\subsection{The screening and validation design}

The two-stage design eliminated half of the work. Stage~A screened all sixteen configurations
under all three policies on the first replication: \(16 \times 3 \times 1 = 48\) simulation
runs. Four configurations were retained as finalists, and Stage~B extended those --- and only
those --- to the two remaining replications: \(4 \times 3 \times 2 = 24\) further runs. The
first replication of the finalists is not repeated, since it was already produced during
screening. The design therefore consumed \(48 + 24 = 72\) simulation runs, against the
\(16 \times 3 \times 3 = 144\) that a full three-replication evaluation of every combination
would have required --- a reduction of exactly \SI{50}{\percent}, without changing which
configuration was ultimately recommended. Of the 48 candidate-policy combinations, 36 therefore
carry single-replication screening metrics only and 12 --- four configurations under three
policies --- carry three-replication aggregates; only the latter are eligible to become the
recommendation.

\Cref{tab:future-validated} reports the validated candidates. The pattern is consistent across
all four: FIFO is infeasible everywhere, with urgent attainment between \SI{60.8}{\percent} and
\SI{70.4}{\percent}, while both urgency-aware policies meet the floors in every replication.

\paragraph{How the three-replication statistics should be read.} Three replications is a small
sample, and the statistics computed from it must be described for what they are. The workflow
reports two summary quantities beyond the mean, and both need qualification.

The first is labelled \emph{p90 cost}. It is the sample \SI{90}{\percent} quantile of three
observations, obtained by linear interpolation between the median and the largest of them ---
four-fifths of the way towards the largest --- and it is therefore dominated by the single most
expensive replication. It is a descriptive indication of
the spread observed, not an estimate of a distributional tail: with three draws, no reliable
statement about a \SI{90}{\percent} exceedance level is available. Read as ``the worst of three
plausible December realisations, approximately'', it is informative; read as a tail-risk figure
it is not, and it is not used as one anywhere in this thesis.

The second is labelled \emph{probability of meeting the SLA targets} in the framework's
interface. It is the share of the three evaluated replications in which the configuration met
both floors, so it can only take the values \(0\), \(\tfrac{1}{3}\), \(\tfrac{2}{3}\) and
\(1\). It is a replication success rate over three trials, not an estimated probability, and no
confidence statement attaches to it. The framework additionally marks a configuration feasible
when that share is at least one half --- a majority rule rather than a requirement that all
three replications succeed. In this run the distinction has no effect on the reported results,
because every validated configuration scored either \(0\) or \(1\): the twelve validated rows
divide into four FIFO rows that failed in all three replications and eight urgency-aware rows
that succeeded in all three.

\begin{table}[htbp]
\centering
\caption[Validated December candidates]{The four configurations validated across three
replications. Costs are means across replications under the forecast run's economic assumptions;
``p90'' is the sample \SI{90}{\percent} quantile of three observations and is descriptive of the
spread, not an estimated tail. ``Feasible in all reps'' records that every one of the three
replications met both floors.}
\label{tab:future-validated}
\small
\begin{tabular}{llrrrrrl}
\toprule
\textbf{Config.} & \textbf{Policy} & \textbf{FTE} & \textbf{Urgent} & \textbf{Normal} &
\textbf{Mean cost} & \textbf{p90 cost} & \textbf{Feasible} \\
 & & & \textbf{(\%)} & \textbf{(\%)} & \textbf{(EUR)} & \textbf{(EUR)} & \textbf{in all reps} \\
\midrule
\multirow{3}{*}{\rgm{s25\_14\_7}} & FIFO & \multirow{3}{*}{46}
   & 60.8 & 94.0 & \num{212595} & \num{225353} & no \\
 & UF   & & 100.0 & 90.2 & \num{163262} & \num{172993} & yes \\
 & RL-3 & & 99.9 & 91.9 & \textbf{\num{158103}} & \num{168923} & yes \\
\midrule
\multirow{3}{*}{\rgm{s26\_14\_7}} & FIFO & \multirow{3}{*}{47}
   & 68.3 & 94.7 & \num{201963} & \num{219823} & no \\
 & UF   & & 100.0 & 91.8 & \num{161162} & \num{171617} & yes \\
 & RL-3 & & 99.9 & 92.7 & \num{158395} & \num{168414} & yes \\
\midrule
\multirow{3}{*}{\rgm{s27\_14\_7}} & FIFO & \multirow{3}{*}{48}
   & 70.1 & 95.1 & \num{200672} & \num{219850} & no \\
 & UF   & & 100.0 & 92.6 & \num{161658} & \num{172551} & yes \\
 & RL-3 & & 99.9 & 93.2 & \num{159760} & \num{169391} & yes \\
\midrule
\multirow{3}{*}{\rgm{s27\_15\_7}} & FIFO & \multirow{3}{*}{49}
   & 70.4 & 95.2 & \num{202838} & \num{222107} & no \\
 & UF   & & 100.0 & 92.7 & \num{164107} & \num{174805} & yes \\
 & RL-3 & & 99.9 & 93.3 & \num{162252} & \num{171637} & yes \\
\bottomrule
\end{tabular}
\end{table}

Two further observations follow from this table. First, the margin between the two
urgency-aware policies is much narrower here than in the retrospective December case ---
between \eur{1855} and \eur{5159}, roughly one to three per cent --- because at these larger
workforces the normal class is no longer close to its floor and there is correspondingly less
for the learned policy to recover. This is the same regime dependence observed in
\Cref{sec:res-policy}, appearing in the forecast run as well.

Second, adding capacity beyond the recommended configuration does not reduce cost. Moving from
\rgm{s25\_14\_7} to \rgm{s27\_15\_7} adds three FTE and improves normal attainment by
1.4 percentage points under RL-3, but raises expected cost by \eur{4149}. The additional
labour is not repaid by the penalties it avoids --- the same conclusion the adaptive search
reached independently in the retrospective run.

Finally, the gap between the mean and the sample p90 cost indicates the spread observed across
the three replications. For the recommended configuration under RL-3 that gap is \eur{10820},
roughly \SI{6.8}{\percent} of the mean cost; for FIFO at the same configuration it is
\eur{12758}. Subject to the caveats above, a planner can read from this both what the plan is
expected to cost and roughly how much more expensive the least favourable of three simulated
December realisations proved to be.

Two features distinguish this output from the retrospective results. First, it is
\emph{spread-aware}: alongside the mean cost, the workflow reports the observed upper end of the
three replications and the share of replications in which the floors were met, so a planner can
distinguish a configuration that succeeded in every replication from one that succeeded on
average. \Cref{fig:future} shows that this distinction is real --- FIFO candidates cluster at a
zero success rate even where their mean cost is not the highest. With three replications this is
an indication rather than a quantified reliability estimate.

Second, the policy pattern of the retrospective run reappears under an independent demand
realisation and a different economic parameter set. At the recommended configuration FIFO
achieves \SI{60.8}{\percent} urgent attainment against RL-3's \SI{99.9}{\percent}, and RL-3's
expected cost of \eur{158103} compares with FIFO's \eur{212595}. The stage-differentiated
behaviour of the learned agent is also visible here: at \rgm{s25\_14\_7} it chose the urgent
queue in \SI{26.7}{\percent} of decisions at picking, \SI{73.2}{\percent} at packing and
\SI{100}{\percent} at dispatch.

\subsection{Stage-differentiated behaviour of the learned policy}
\label{sec:res-rl-behaviour}

The forecast run records how often the learned agent chose each class at each stage, which gives
a direct view of its behaviour rather than only of its outcomes. At the recommended
configuration the urgent queue was selected in \SI{26.7}{\percent} of decisions at picking,
\SI{73.2}{\percent} at packing and \SI{100}{\percent} at dispatch.

The spread is striking because a single set of network parameters produced all three figures.
The agent is not three policies; it is one policy conditioned on, among other things, a feature
identifying which stage the decision is being taken at. It has evidently learned to apply very
different logic depending on where in the flow it is acting --- close to indifference at
picking, a strong urgent bias at packing, and absolute priority at dispatch.

A plausible reading is that the stages differ in how consequential a deferral is. An urgent
order deferred at picking still has the whole of its remaining path in which to recover, and the
picking queue is long, so the opportunity cost of serving it ahead of many normal orders is
high. At dispatch, by contrast, the order is one operation from completion and any deferral
directly translates into lateness. This is an interpretation of the observed behaviour, not a
demonstrated mechanism: nothing in this study isolates why the agent behaves this way, and it
was not designed to do so. What can be stated is that a stage-independent rule cannot express
this pattern at all, because it applies identical logic everywhere.

\subsection{Comparability between the two runs}

It must be emphasised that this run used different economic assumptions from the retrospective
run --- a late urgent order is penalised at \eur{15} rather than \eur{20}, a late normal order
at \eur{10} rather than \eur{5}, and labour costs \eur{18} per hour rather than \eur{15}. Its
monetary values are therefore not comparable with those in
\Cref{tab:recommendations,tab:cases,tab:cost-split}. Notably, the two runs weight the two
classes of failure quite differently: the retrospective run penalises an urgent failure four
times as heavily as a normal one, while the forecast run penalises it only \num{1.5} times as
heavily. A policy that sacrifices normal orders to protect urgent ones is therefore
comparatively more attractive under the retrospective assumptions than under the forecast ones.

The comparison that \emph{is} valid is the one made within each run, between candidates and
policies evaluated under a common parameter set. That the same qualitative pattern --- FIFO
infeasible, the two urgency-aware policies close, RL-3 ahead where the normal class binds ---
appears under both parameter sets is mild evidence that the pattern is not an artefact of one
particular choice of penalties.

\section{The annual view}
\label{sec:res-annual}

The configurations examined in detail above were selected because they illustrate distinct
operating regimes. To guard against the impression that the argument rests on them, this section
reports the year as a whole.

\subsection{Three seasons, three different problems}

Read across all twelve months, the results separate into three regimes that pose different
planning problems.

\paragraph{The summer trough (May to August).} Demand falls to between \num{8400} and
\num{10800} orders per month and the urgent share to \SIrange{7.9}{9.0}{\percent}. Recommended
workforces are 7 to 8 FTE, penalties at the recommended configuration are trivial (\eur{155} to
\eur{1495}, against labour costs of \eur{16800} to \eur{19200}), and the policies differ by tens
of euros. FIFO remains feasible
at 8 to 11 of the 16 candidates. In this regime the planning question is simply how far capacity
can be reduced without breaching the floors; the sequencing rule is close to irrelevant.

\paragraph{The shoulder months (March, April, September, October).} Demand runs between
\num{16800} and \num{20400} orders with urgent shares of \SIrange{12.1}{15.2}{\percent}.
Recommended workforces are 13 to 17 FTE. FIFO is now feasible at only 11 to 14 candidates and
begins to fail at the leaner configurations, while the two urgency-aware policies remain feasible
at 14 to 16. Urgent-First is selected as lowest-cost in all four of these months, and its margin
over RL-3 is small --- between \eur{20} and \eur{85}. This is the regime in which an urgency-aware
rule has become necessary but a sophisticated one has not.

\paragraph{The campaign months (January, February, November, December).} Demand rises to between
\num{26400} and \num{40800} orders with urgent shares of \SIrange{18.0}{25.4}{\percent}, and
recommended workforces to 25 to 38 FTE. FIFO is feasible at \emph{none} of the 16 candidates in
any of the four --- the only months in which that happens, and the reason the candidate range
itself, rather than any property of capacity in general, bounds the claim. Penalty costs become the dominant cost component, and the spread between
policies runs into five figures. RL-3 is selected in all four. This is the regime in which the
sequencing decision determines whether a viable plan exists at all.

\subsection{The recommendation is not a smooth function of volume}

\Cref{tab:recommendations} shows that the recommended workforce does not scale proportionally
with demand. December handles \num{4.86} times July's volume but requires \num{5.43} times the
workforce, whereas March handles \num{1.56} times May's volume with \num{1.63} times the
workforce. Two effects pull in opposite directions. Workload per order rises with volume, as
\Cref{sec:res-demand} showed, so peak months need proportionally more capacity per order ---
which is why December's ratio exceeds its volume ratio. Against that, integer rounding at three
stages is proportionally more costly in small configurations, which inflates the summer months'
workforces relative to a smooth curve and compresses the ratios between them.

The practical consequence is that a planner cannot scale a known good configuration by a demand
ratio. The recommendation for a month with twice the volume is not twice the workforce, and the
error runs in both directions depending on where in the seasonal cycle the reference month sits.

\subsection{Consistency of the qualitative findings}

Three findings hold in every month examined, without exception:

\begin{itemize}
  \item picking is the primary pressure location (12 of 12);
  \item the analytical estimate exceeds the simulated recommendation (12 of 12);
  \item where the adaptive search ran, the additional FTE was rejected (9 of 9).
\end{itemize}

\noindent Two hold with clear exceptions that are themselves explicable:

\begin{itemize}
  \item FIFO is never selected as lowest-cost (12 of 12), but it \emph{is} feasible in the
        eight non-campaign months, so it is not universally unusable --- only unusable under
        pressure;
  \item RL-3's normal-class advantage over Urgent-First holds throughout the feasible region
        (167 configurations, worst case \(-0.08\) points) and reverses entirely outside it
        (25 configurations, mean \(-17.69\) points).
\end{itemize}

\section{Summary of observations}
\label{sec:res-summary}

\begin{enumerate}
  \item Orders impose markedly different workload at the three stages
        (\(1 : 0.72 : 0.34\)), and seasonal growth concentrates on picking.
  \item The analytical capacity estimate exceeded the simulated recommendation in all twelve
        months, by up to nine FTE in December.
  \item FIFO was feasible in \num{84} of \num{192} runs against \num{166} and \num{167} for the
        two urgency-aware policies. Its infeasible configurations occur in every month, but
        \emph{complete} monthly infeasibility --- no feasible configuration among the sixteen
        --- occurs only in the four campaign months.
  \item Policy differences are large under capacity pressure and negligible under slack: at
        October \rgm{s11\_7\_5} the three policies differ in total cost by \eur{30}.
  \item At the December recommended configuration, RL-3 was the only feasible policy, at
        \SI{13.2}{\percent} lower cost than Urgent-First.
  \item Picking was the primary pressure location in all twelve months and accounted for
        \SIrange{83.7}{100}{\percent} of late orders' accumulated waiting.
  \item The adaptive search tested one additional picking FTE in nine months and rejected it in
        all nine.
  \item The forecast workflow reproduced the qualitative policy pattern under independently
        generated demand, reporting results across three replications rather than a single
        point estimate. Three replications describe the spread observed; they do not quantify
        risk.
\end{enumerate}
